بازی شطرنج، به‌عنوان یکی از مسائل شناخته‌شده در حوزه هوش مصنوعی، با پیشرفت موتورهای شطرنجی به سطحی از دقت رسیده است که تحلیل تاکتیکی و ارزیابی موقعیت‌ها را با دقت بسیار بالا امکان‌پذیر می‌کند. با این حال، خروجی موتورهایی مانند \lr{Stockfish} معمولاً به‌صورت ارزیابی عددی و خطوط اصلی بازی ارائه می‌شود و برای بازیکنان، به‌ویژه بازیکنان مبتدی و متوسط، لزوماً قابل فهم و آموزشی نیست. از سوی دیگر، مدل‌های زبانی بزرگ توانایی بالایی در تولید متن روان دارند، اما در توضیح موقعیت‌های شطرنجی ممکن است دچار خطا یا ارائه تحلیل‌های نادرست شوند. هدف این پروژه، طراحی و پیاده‌سازی سامانه‌ای برای تولید توضیحات آموزشی و قابل فهم درباره حرکات شطرنج است که ضمن بهره‌گیری از توانایی تولید زبان طبیعی، به تحلیل دقیق موتور شطرنج نیز متکی باشد.

برای ارزیابی منصفانه و یکسان مدل‌ها، یک مجموعه‌داده آزمون متوازن شامل ۶۰نمونه از موقعیت‌ها و حرکات مختلف شطرنج از بازی‌های واقعی تهیه و به‌صورت دستی بررسی و برچسب‌گذاری شد. این مجموعه شامل ۲۰نمونه از هر یک از سطوح مبتدی، متوسط و پیشرفته است و برای هر نمونه، وضعیت صفحه قبل و بعد از حرکت، اطلاعات حرکت، تحلیل‌های مرتبط و یک توضیح مرجع آموزشی در نظر گرفته شده است. سپس ۱۱مدل زبانی با اندازه‌ها و قابلیت‌های متفاوت، از مدل کوچک مبتنی بر \lr{Qwen3-1.7B} با \lr{۱٫۷} میلیارد پارامتر تا مدل‌های جدیدتری مانند \lr{GPT-5.5} و \lr{Qwen3.7}بر روی این مجموعه آزمایش شدند. خروجی مدل‌ها نیز توسط یک مدل زبانی مستقل، \lr{Gemma 4 12B}، بر اساس معیارهایی شامل صحت، پیروی از دستور، مفیدبودن آموزشی، میزان انطباق با توضیح مرجع و روانی متن ارزیابی شد.

از نظر نرم‌افزاری، سامانه به‌صورت یک پایپالین چندمرحله‌ای طراحی شده است که در آن ابتدا اطلاعات موقعیت و حرکت شطرنج دریافت شده و با استفاده از \lr{Stockfish}تحلیل می‌شود؛ سپس ویژگی‌ها و تغییرات مهم صفحه توسط ماژول پردازش شطرنج استخراج شده و در قالب یک ورودی ساختاریافته به مدل زبانی داده می‌شود. این ساختار امکان آزمایش مدل‌های مختلف و مقایسه خروجی آن‌ها را در یک محیط یکسان فراهم می‌کند.

مدل اختصاصی مبتنی بر \lr{Qwen3-1.7B} با \lr{۱٫۷} میلیارد پارامتر، امتیاز کلی \lr{۲٫۸۶} از ده را کسب کرد. این امتیاز تقریباً برابر با مدل ۳ میلیاردپارامتری \lr{Llama 3.2} و نزدیک به \lr{Gemma 3 12B} با \lr{۲٫۹۵} و \lr{Gemma 4 26B A4B} با \lr{۲٫۹۰} است. بنابراین، مدل \lr{۱٫۷} میلیاردپارامتری در این مجموعه توانسته است میانگین امتیازی نزدیک به چند مدل بزرگ‌تر داشته باشد؛ هرچند این نتیجه به معنی برابری عمومی توانایی‌ها یا اثبات اثر آموزش دامنه نیست. بهترین مدل آزموده‌شده امتیاز \lr{۷٫۳۱} دارد و شکاف روانی و صحت در مدل کوچک همچنان باقی است. در مجموع، پروژه چارچوبی مشترک برای تولید و ارزیابی توضیحات شطرنجی فراهم می‌کند و ظرفیت مدل‌های سبک را در کنار ضرورت بهبود صحت محتوایی نشان می‌دهد.

واژگان کلیدی: هوش مصنوعی شطرنج، مدل‌های زبانی، تولید زبان طبیعی، ارزیابی مدل‌های زبانی، توضیح‌پذیری، تحلیل شطرنج